\chapter{Gaussian calculus on self-shrinkers}\label{sec:shrinker-calculus}
We develop the weighted identities behind the variation theory of self-shrinkers. The primary reference is \cite{CM12}, Section 3.3, pp.~772--775. All integrals on immersed hypersurfaces are taken on the parameter manifold with its induced area measure, so multiplicity is retained. Write $[u]=\int_\Sigma u\rho\,dA$, where $\rho=(4\pi)^{-n/2}e^{-|x|^2/4}$. Our scalar convention is $\vec H=-H\nu$ and $H=\operatorname{div}_\Sigma\nu$.

\begin{definition}[Gaussian drift and polynomial area growth]\label{def:gaussian-drift}
On a smooth immersed $n$-hypersurface $\Sigma\subset\mathbb R^{n+1}$, $n\geq1$, define
\[
 \mathcal D u=\Delta_\Sigma u-\tfrac12\langle x^T,\nabla_\Sigma u\rangle.
\]
Polynomial area growth means that $\operatorname{Area}(\Sigma\cap B_R(0))\leq C(1+R)^d$ for all $R\geq1$, for some finite $C,d$. The intersection here means the inverse image under the immersion. This condition controls Gaussian tails; properness separately makes ambient radial cutoffs compactly supported. See \cite{CM12}, Lemma 3.8 and the hypotheses of Lemma 3.25.
\end{definition}

\begin{lemma}[Weighted integration by parts with compact support]\label{lem:gaussian-integration-parts}
On a smooth boundaryless immersed hypersurface, if $u\in C_c^1(\Sigma)$ and $v\in C^2(\Sigma)$, then
\[
 \mathcal Dv=\rho^{-1}\operatorname{div}_\Sigma(\rho\nabla v),
 \qquad [u\mathcal Dv]=-[\langle\nabla u,\nabla v\rangle].
\]
Neither the shrinker equation nor finite total Gaussian area is needed.
\end{lemma}
\begin{proof}
The tangential derivative of $\log\rho$ is $-x^T/2$. Expand the divergence using the product rule. Integrate $\operatorname{div}_\Sigma(u\rho\nabla v)$; it has compact support and its integral is zero by the divergence theorem. This proves the identity, which is \cite{CM12}, Lemma 3.8.
\end{proof}

\begin{lemma}[Global weighted integration by parts]\label{lem:gaussian-global-parts}
Suppose $\Sigma$ is a connected complete smooth immersed hypersurface without boundary, and $u,v\in C^2(\Sigma)$ satisfy
\[
 [\,|u\nabla v|+|\nabla u||\nabla v|+|u\mathcal Dv|\,]<\infty.
\]
Then $[u\mathcal Dv]=-[\langle\nabla u,\nabla v\rangle]$. If these conditions also hold with $u,v$ interchanged, then $[u\mathcal Dv]=[v\mathcal Du]$.
\end{lemma}
\begin{proof}
Take intrinsic distance cutoffs $0\leq\eta_j\leq1$, equal to one on the intrinsic ball of radius $j$ about a fixed point, zero outside radius $j+1$, and with $|\nabla\eta_j|\leq1$ almost everywhere. Completeness makes their supports compact. The compact-support identity extends to these Lipschitz tests by local smooth approximation. Apply it to $\eta_ju$ and $v$:
\[
 [\eta_ju\mathcal Dv]=-[\eta_j\langle\nabla u,\nabla v\rangle]
 -[u\langle\nabla\eta_j,\nabla v\rangle].
\]
Dominated convergence applies to the first two terms. The last tends to zero since its absolute value is at most the tail integral of $|u\nabla v|\rho$. Interchanging the functions gives symmetry. These are the integrability requirements in \cite{CM12}, Corollary 3.10; completeness alone does not remove them.
\end{proof}

\begin{lemma}[Gaussian moments from area growth]\label{lem:gaussian-polynomial-moments}
If a smooth immersed hypersurface has polynomial area growth, then for every $a>0$ and $q\geq0$,
\[
 \int_\Sigma(1+|x|)^q e^{-a|x|^2}\,dA<\infty.
\]
The integral over $\{|x|>R\}$ tends to zero as $R\to\infty$.
\end{lemma}
\begin{proof}
The part inside $B_1$ has finite area by hypothesis. On $B_{j+1}\setminus B_j$, the integral is bounded by $C(j+2)^{q+d}e^{-aj^2}$. This series converges: the quadratic exponential eventually dominates any polynomial. Its tails tend to zero. This also justifies differentiating Gaussian integrals in center and positive scale on compact parameter sets, since each derivative is bounded by a polynomial times $e^{-a|x|^2}$ for some $a>0$.
\end{proof}

\begin{proposition}[Coordinate and radial drift identities]\label{prop:shrinker-drift-identities}
On any smooth immersed self-shrinker $\Sigma^n\subset\mathbb R^{n+1}$,
\[
 \mathcal D x_i=-\tfrac12x_i,\qquad
 \mathcal D|x|^2=2n-|x|^2,\qquad
 |\nabla |x|^2|^2=4|x^T|^2.
\]
These are pointwise identities and require no growth hypothesis.
\end{proposition}
\begin{proof}
The immersion satisfies $\Delta_\Sigma x=\vec H=-x^\perp/2$. The drift subtracts $x^T/2$, giving $\mathcal Dx=-x/2$. The product rule for $\mathcal D$ and $\sum_i|\nabla x_i|^2=n$ give the radial identity. Finally $\nabla |x|^2=2x^T$. Compare \cite{CM12}, Lemma 3.20.
\end{proof}

\begin{proposition}[Gaussian center, radius, and mixed moments]\label{prop:shrinker-basic-moments}
Let $\Sigma^n\subset\mathbb R^{n+1}$ be a connected complete smooth immersed self-shrinker without boundary and with polynomial area growth. Then
\[
 [x]=0,\qquad [|x|^2]=2n[1],\qquad [x|x|^2]=0.
\]
For every constant vector $v\in\mathbb R^{n+1}$,
\[
 [\langle x,v\rangle^2]=2[|v^T|^2].
\]
\end{proposition}
\begin{proof}
All products of the coordinate functions, $|x|^2$, their first derivatives, and their drift derivatives that occur here are polynomially bounded. Lemma \ref{lem:gaussian-polynomial-moments} therefore verifies each hypothesis of Lemma \ref{lem:gaussian-global-parts}. Apply that lemma with $u=1$ and $v=x_i$ or $|x|^2$ to obtain the first two identities from Proposition \ref{prop:shrinker-drift-identities}. Symmetry with $u=x_i$, $v=|x|^2$ gives
\[
 -\tfrac12[x_i|x|^2]=[|x|^2\mathcal Dx_i]
 =[x_i\mathcal D|x|^2]=-[x_i|x|^2],
\]
so this mixed moment is zero. Finally, setting $u=v=\langle x,v\rangle$ in integration by parts gives $-\tfrac12[\langle x,v\rangle^2]=-[|v^T|^2]$. These are the identities of \cite{CM12}, Lemma 3.25.
\end{proof}

\begin{proposition}[The fourth moment and the scale variance]\label{prop:shrinker-fourth-moment}
Let $\Sigma$ satisfy the hypotheses of Proposition \ref{prop:shrinker-basic-moments} and be two-sided, with $H=\langle x,\nu\rangle/2$. Then
\[
 [|x|^4]=4n(n+2)[1]-16[H^2],
 \qquad
 \left[\left(\tfrac{|x|^2}{4}-\tfrac n2\right)^2\right]
 =\tfrac n2[1]-[H^2].
\]
In particular $[H^2]\leq n[1]/2$.
\end{proposition}
\begin{proof}
Integration by parts with $u=v=|x|^2$ is justified as above and gives
\[
 2n[|x|^2]-[|x|^4]=-4[|x^T|^2]
 =-4[|x|^2]+16[H^2].
\]
Insert $[|x|^2]=2n[1]$ to obtain the first formula. Expand the square and substitute both moment identities to get the second. Its left side is nonnegative. This is \cite{CM12}, Lemma 3.25 and Corollary 3.34.
\end{proof}

\begin{corollary}[A radial constraint for closed shrinkers]\label{cor:shrinker-radius-constraint}
A connected smooth closed embedded self-shrinker $\Sigma^n\subset\mathbb R^{n+1}$ satisfies
\[
 \min_\Sigma |x|\leq\sqrt{2n}\leq\max_\Sigma |x|.
\]
If it lies entirely in $\overline B_{\sqrt{2n}}$ or entirely in its exterior $\{|x|\geq\sqrt{2n}\}$, it is the round sphere of radius $\sqrt{2n}$ centered at zero.
\end{corollary}
\begin{proof}
The positive weighted measure and $[|x|^2-2n]=0$ force both inequalities. Under either one-sided radial condition, the continuous integrand has one sign and zero integral, so $|x|^2=2n$ everywhere. The inclusion in the sphere is a local diffeomorphism because both manifolds have dimension $n$ and the immersion has full rank. The image is open, and compactness makes it closed. The sphere is connected for $n\geq1$, so the image is the whole sphere. Embeddedness identifies it with the round sphere rather than a covering parametrization.
\end{proof}
